\section{Proofs for Section 4}

This appendix collects the algebra omitted from Section \ref{sec:hub_efficiency}.

\subsection{Preliminaries}

Recall from Section \ref{sec:model} that
\[
P_A=\frac{N_A(N_A-1)}{2},
\qquad
Q_{AB}=N_A N_B,
\qquad
m_A(h)=m_A^0+\bigl(m_A^1-m_A^0\bigr)h,
\]
with
\[
0\le m_A^0<m_A^1\le 1.
\]
In the benchmark case without pipeline policy ($b=0$), welfare is
\begin{equation}
W(0,0)
=
v_L P_A m_A^0
+
\beta\Bigl\{
v_L P_A m_A^0\bigl(1-\delta m_A^0\bigr)
+
v_X Q_{AB}q_0
\Bigr\},
\label{eq:app4_W00}
\end{equation}
and
\begin{equation}
W(1,0)
=
v_L P_A m_A^1
+
\beta\Bigl\{
v_L P_A m_A^1\bigl(1-\delta m_A^1\bigr)
+
v_X Q_{AB}q_0
\Bigr\}
-F.
\label{eq:app4_W10}
\end{equation}
Hence the net welfare gain from a hub alone is
\begin{align}
\Delta_H
&:=W(1,0)-W(0,0)
\nonumber\\
&=
v_L P_A\Bigl[
m_A^1-m_A^0
+
\beta\bigl(
m_A^1(1-\delta m_A^1)-m_A^0(1-\delta m_A^0)
\bigr)
\Bigr]-F.
\label{eq:app4_Delta_start}
\end{align}
The cross-regional baseline term cancels exactly:
\[
\beta v_XQ_{AB}q_0-\beta v_XQ_{AB}q_0=0.
\]
Thus the hub-alone comparison depends only on local primitives.

Now,
\begin{align}
m_A^1(1-\delta m_A^1)-m_A^0(1-\delta m_A^0)
&=
(m_A^1-m_A^0)-\delta\bigl((m_A^1)^2-(m_A^0)^2\bigr)
\nonumber\\
&=
(m_A^1-m_A^0)\bigl[1-\delta(m_A^1+m_A^0)\bigr].
\label{eq:app4_local_diff}
\end{align}
Substituting \eqref{eq:app4_local_diff} into \eqref{eq:app4_Delta_start} yields
\begin{align}
\Delta_H
&=
v_L P_A(m_A^1-m_A^0)
\Bigl[(1+\beta)-\beta\delta(m_A^1+m_A^0)\Bigr]-F.
\label{eq:app4_Delta_factorized}
\end{align}
Define
\begin{equation}
\Lambda_A:=(1+\beta)-\beta\delta(m_A^1+m_A^0),
\label{eq:app4_Lambda}
\end{equation}
so that
\begin{equation}
\Delta_H=v_L P_A(m_A^1-m_A^0)\Lambda_A-F.
\label{eq:app4_Delta_Lambda}
\end{equation}

\subsection{Proof of Proposition \ref{prop:hub_efficiency}}

We prove each part in turn.

\paragraph{Part 1: the case $\Lambda_A\le 0$.}
Since $v_L>0$, $P_A>0$ for all $N_A\ge 2$, and $m_A^1-m_A^0>0$, the product
\[
v_L P_A(m_A^1-m_A^0)\Lambda_A
\]
is weakly negative whenever $\Lambda_A\le 0$. Therefore, from \eqref{eq:app4_Delta_Lambda},
\[
\Delta_H
=
v_L P_A(m_A^1-m_A^0)\Lambda_A-F
<0
\]
for every $F>0$. Hence the hub is never socially efficient.

\paragraph{Part 2: the case $\Lambda_A>0$.}
Under $\Lambda_A>0$, the condition $\Delta_H\ge 0$ is equivalent to
\[
v_L P_A(m_A^1-m_A^0)\Lambda_A\ge F.
\]
Since all terms except $P_A$ are strictly positive, this can be rearranged as
\begin{equation}
P_A
\ge
\frac{F}{v_L(m_A^1-m_A^0)\Lambda_A}.
\label{eq:app4_pair_threshold}
\end{equation}
This proves \eqref{eq:pair_threshold}.

Using $P_A=N_A(N_A-1)/2$, \eqref{eq:app4_pair_threshold} becomes
\[
\frac{N_A(N_A-1)}{2}
\ge
\frac{F}{v_L(m_A^1-m_A^0)\Lambda_A},
\]
or equivalently
\[
N_A^2-N_A-\frac{2F}{v_L(m_A^1-m_A^0)\Lambda_A}\ge 0.
\]
The positive root of the corresponding quadratic equation is
\[
\frac{1}{2}\left(
1+\sqrt{1+\frac{8F}{v_L(m_A^1-m_A^0)\Lambda_A}}
\right).
\]
Hence, treating $N_A$ as continuous, the hub is efficient if and only if
\begin{equation}
N_A\ge
N_H^*
:=
\frac{1}{2}\left(
1+\sqrt{1+\frac{8F}{v_L(m_A^1-m_A^0)\Lambda_A}}
\right),
\label{eq:app4_Nstar}
\end{equation}
which is exactly \eqref{eq:Nstar}. For integer $N_A$, replace $N_H^*$ by its ceiling. This completes the proof.
\qed

\subsection{Derivation of the Equivalent Threshold Characterizations}

Starting from \eqref{eq:app4_Delta_Lambda}, the condition $\Delta_H\ge 0$ is equivalent to
\[
F\le v_L P_A(m_A^1-m_A^0)\Lambda_A.
\]
Thus the critical fixed cost is
\begin{equation}
F_H^*(N_A):=v_L P_A(m_A^1-m_A^0)\Lambda_A,
\label{eq:app4_Fstar}
\end{equation}
which is \eqref{eq:Fstar}.

Next, solve for the critical redundancy rate. Rewrite $\Delta_H\ge 0$ as
\[
v_L P_A(m_A^1-m_A^0)\Bigl[(1+\beta)-\beta\delta(m_A^1+m_A^0)\Bigr]\ge F.
\]
Dividing by the positive term $v_L P_A(m_A^1-m_A^0)$ gives
\[
(1+\beta)-\beta\delta(m_A^1+m_A^0)
\ge
\frac{F}{v_L P_A(m_A^1-m_A^0)}.
\]
Hence
\[
\beta\delta(m_A^1+m_A^0)
\le
(1+\beta)-\frac{F}{v_L P_A(m_A^1-m_A^0)},
\]
so
\begin{align}
\delta
&\le
\frac{1+\beta}{\beta(m_A^1+m_A^0)}
-
\frac{F}{\beta v_L P_A(m_A^1-m_A^0)(m_A^1+m_A^0)}
\nonumber\\
&=:\delta_H^*(N_A),
\label{eq:app4_deltastar}
\end{align}
which is \eqref{eq:deltastar}.

\subsection{Proof of Proposition \ref{prop:comparative_statics_hub}}

Throughout this proof we use
\[
\Delta_H=v_L P_A(m_A^1-m_A^0)\Lambda_A-F,
\qquad
\Lambda_A=(1+\beta)-\beta\delta(m_A^1+m_A^0).
\]

\paragraph{Derivative with respect to $N_A$.}
Treating $N_A$ as continuous,
\[
P_A=\frac{N_A(N_A-1)}{2}
\quad\Rightarrow\quad
\frac{\partial P_A}{\partial N_A}=\frac{2N_A-1}{2}.
\]
Hence
\begin{align}
\frac{\partial \Delta_H}{\partial N_A}
&=
v_L\frac{\partial P_A}{\partial N_A}(m_A^1-m_A^0)\Lambda_A
\nonumber\\
&=
\frac{v_L}{2}(2N_A-1)(m_A^1-m_A^0)\Lambda_A,
\label{eq:app4_dDelta_dNA}
\end{align}
which is \eqref{eq:dDelta_dN}. If $\Lambda_A>0$, then $\partial\Delta_H/\partial N_A>0$ because $v_L>0$, $2N_A-1>0$, and $m_A^1-m_A^0>0$.

For integer $N_A$, the same monotonicity follows from forward differences:
\begin{align}
\Delta_H(N_A+1)-\Delta_H(N_A)
&=
v_L\bigl(P_A(N_A+1)-P_A(N_A)\bigr)(m_A^1-m_A^0)\Lambda_A
\nonumber\\
&=
v_LN_A(m_A^1-m_A^0)\Lambda_A
>0
\quad\text{if }\Lambda_A>0.
\label{eq:app4_forward_difference}
\end{align}

\paragraph{Derivative with respect to $F$.}
From \eqref{eq:app4_Delta_Lambda},
\begin{equation}
\frac{\partial \Delta_H}{\partial F}=-1<0,
\label{eq:app4_dDelta_dF}
\end{equation}
which is \eqref{eq:dDelta_dF}.

\paragraph{Derivative with respect to $\delta$.}
Differentiating $\Lambda_A=(1+\beta)-\beta\delta(m_A^1+m_A^0)$ gives
\[
\frac{\partial \Lambda_A}{\partial \delta}=-\beta(m_A^1+m_A^0).
\]
Therefore
\begin{align}
\frac{\partial \Delta_H}{\partial \delta}
&=
v_L P_A(m_A^1-m_A^0)\frac{\partial \Lambda_A}{\partial \delta}
\nonumber\\
&=
-\beta v_L P_A(m_A^1-m_A^0)(m_A^1+m_A^0)
<0,
\label{eq:app4_dDelta_ddelta}
\end{align}
which is \eqref{eq:dDelta_ddelta}.

\paragraph{Derivative with respect to $m_A^1$.}
Starting from
\[
\Delta_H=v_L P_A(m_A^1-m_A^0)\Bigl[(1+\beta)-\beta\delta(m_A^1+m_A^0)\Bigr]-F,
\]
differentiate with respect to $m_A^1$:
\begin{align}
\frac{\partial \Delta_H}{\partial m_A^1}
&=
v_L P_A
\left(
\Bigl[(1+\beta)-\beta\delta(m_A^1+m_A^0)\Bigr]
+
(m_A^1-m_A^0)(-\beta\delta)
\right)
\nonumber\\
&=
v_L P_A\Bigl[(1+\beta)-2\beta\delta m_A^1\Bigr].
\label{eq:app4_dDelta_dm1_raw}
\end{align}
This is exactly \eqref{eq:dDelta_dm1}.

To verify the sign, note that the regularity condition \eqref{eq:regularity} implies
\[
1-\delta m_A^1>0
\quad\Longrightarrow\quad
\delta m_A^1<1.
\]
Hence
\[
2\beta\delta m_A^1<2\beta.
\]
Since $\beta\in(0,1]$, we have $1+\beta\ge 2\beta$, with equality only when $\beta=1$. Therefore,
\[
(1+\beta)-2\beta\delta m_A^1>(1+\beta)-2\beta=1-\beta\ge 0,
\]
where the strict inequality follows from $\delta m_A^1<1$. Thus
\begin{equation}
\frac{\partial \Delta_H}{\partial m_A^1}>0.
\label{eq:app4_dDelta_dm1_sign}
\end{equation}

\paragraph{Derivative with respect to $m_A^0$.}
Similarly,
\begin{align}
\frac{\partial \Delta_H}{\partial m_A^0}
&=
v_L P_A
\left(
-\Bigl[(1+\beta)-\beta\delta(m_A^1+m_A^0)\Bigr]
+
(m_A^1-m_A^0)(-\beta\delta)
\right)
\nonumber\\
&=
-\,v_L P_A\Bigl[(1+\beta)-2\beta\delta m_A^0\Bigr].
\label{eq:app4_dDelta_dm0_raw}
\end{align}
This is exactly \eqref{eq:dDelta_dm0}.

Because $m_A^0<m_A^1$ and $\delta m_A^1<1$, we also have $\delta m_A^0<1$. Repeating the same sign argument as above gives
\[
(1+\beta)-2\beta\delta m_A^0>0,
\]
so
\begin{equation}
\frac{\partial \Delta_H}{\partial m_A^0}<0.
\label{eq:app4_dDelta_dm0_sign}
\end{equation}

This proves Proposition \ref{prop:comparative_statics_hub}.
\qed
